
\chapter{مفاهیم اولیه}
\label{ch:prelim}

این فصل مفاهیم و ابزارهایی را معرفی می‌کند که در فصل‌های بعدی به کار می‌روند.
ترتیب مطالب از سیگنال گفتار و ویژگی‌های آوایی آغاز می‌شود، به مؤلفه‌های زنجیره
آبشاری و سپس به معماری مدل‌های سرتاسری گفتار‌به‌گفتار می‌رسد، و با سنجه‌ها و
مفاهیم آماری ارزیابی پایان می‌یابد.

\section{سیگنال گفتار و ویژگی‌های آوایی}
\label{sec:features}

گفتار در این پروژه به‌صورت سیگنال تک‌کاناله با بسامد نمونه‌برداری
\N{16} کیلوهرتز و نمونه‌های صحیح ۱۶‌بیتی نمایش داده می‌شود. سیگنال پیوسته
$x[n]$ به «فریم»هایی با طول ثابت تقسیم می‌شود؛ در این پروژه طول فریم
\N{20} میلی‌ثانیه است، که با \N{320} نمونه معادل است. تقسیم به فریم بر این
فرض استوار است که خواص آماری گفتار در بازه‌های کوتاه تقریباً ثابت‌اند.

\subsection{انرژی و نرخ گذر از صفر}

ساده‌ترین ویژگی، \مهم{انرژی فریم} است. برای فریم $t$ با $N$ نمونه، انرژی
لگاریتمی به‌صورت زیر محاسبه می‌شود:

\begin{equation}
E_t = \log\!\left(\frac{1}{N}\sum_{n=1}^{N} x_t^2[n] + \epsilon\right),
\label{eq:energy}
\end{equation}

که در آن $\epsilon$ یک مقدار مثبت کوچک برای پرهیز از لگاریتم صفر است. انرژی
به‌تنهایی برای تشخیص گفتار کافی نیست، زیرا نویز پرتوان نیز انرژی بالایی
دارد. ویژگی مکمل، \مهم{نرخ گذر از صفر} است که تعداد تغییر علامت سیگنال را
در فریم می‌شمارد:

\begin{equation}
Z_t = \frac{1}{2(N-1)}\sum_{n=2}^{N}
      \bigl|\operatorname{sgn}(x_t[n]) - \operatorname{sgn}(x_t[n-1])\bigr|.
\label{eq:zcr}
\end{equation}

آواهای واک‌دار نرخ گذر از صفر پایین و آواهای سایشی نرخ بالا دارند؛ ترکیب این
دو ویژگی مبنای آشکارسازهای کلاسیک تشخیص فعالیت گفتاری
است~\cite{sohn1999vad}.

\subsection{بسامد پایه و روش \en{YIN}}

\مهم{بسامد پایه} یا $F_0$، بسامد ارتعاش تارهای صوتی و همان چیزی است که به‌طور
ادراکی به‌عنوان زیر‌و‌بمی شنیده می‌شود. تخمین $F_0$ در این پروژه با روش
\en{YIN}~\cite{decheveigne2002yin} انجام می‌شود. این روش بر تابع اختلاف
مربعی هم‌بستگی استوار است:

\begin{equation}
d_t(\tau) = \sum_{n=1}^{W} \bigl(x_t[n] - x_t[n+\tau]\bigr)^2,
\label{eq:yin-diff}
\end{equation}

که در آن $\tau$ تأخیر و $W$ طول پنجره تحلیل است. سپس نسخه «هنجارشده
تجمعی» این تابع ساخته می‌شود و نخستین کمینه محلی که از یک آستانه پایین‌تر
بیفتد به‌عنوان دوره تناوب $\hat{\tau}$ انتخاب می‌گردد؛ بسامد پایه برابر
$F_0 = f_s / \hat{\tau}$ خواهد بود. اگر هیچ کمینه‌ای از آستانه پایین‌تر
نرود، فریم «بی‌واک» علامت‌گذاری می‌شود. در این پروژه بازه جست‌وجوی
بسامد پایه \N{60} تا \N{400} هرتز در نظر گرفته شده است که گستره گویندگان
مرد و زن را پوشش می‌دهد.

اهمیت $F_0$ در تشخیص وقفه آن است که آغاز یک گفتار انسانی معمولاً با ظهور
ساختار واک‌دار همراه است، در حالی که نویز محیط یا صدای پخش بلندگو الگوی
زیر‌و‌بمی پایدار ندارد.

\subsection{ضرایب کپستروم بسامد مل}

\en{MFCC} یا \مهم{ضرایب کپستروم بسامد مل}~\cite{davis1980mfcc} پرکاربردترین
نمایش کوتاه‌مدت گفتار است و طیف فریم را به شکلی فشرده می‌کند که با ادراک
شنوایی انسان هم‌خوان باشد. مراحل محاسبه چنین است:

\begin{enumerate}
\item اعمال پنجره (مثلاً پنجره هَمینگ) و محاسبه تبدیل فوریه گسسته فریم؛
\item محاسبه توان طیفی و عبور آن از یک بانک فیلتر مثلثی که مراکز آن روی
      مقیاس مل توزیع شده‌اند. رابطه مقیاس مل با بسامد خطی چنین است:
      \begin{equation}
      m(f) = 2595 \log_{10}\!\left(1 + \frac{f}{700}\right);
      \label{eq:mel}
      \end{equation}
\item لگاریتم‌گیری از خروجی هر فیلتر؛
\item اعمال تبدیل کسینوسی گسسته برای ناهم‌بسته‌کردن مؤلفه‌ها و نگه‌داشتن
      نخستین ضرایب.
\end{enumerate}

در این پروژه \N{13} ضریب نگه داشته می‌شود. برای آن‌که ویژگی‌ها اطلاعات
«تغییر» در زمان را نیز حمل کنند، مشتق مرتبه اول این ضرایب، که به
«دلتا» شناخته می‌شود، نیز محاسبه و به بردار ویژگی افزوده می‌گردد. مجموع
ویژگی‌های هر فریم بنابراین شامل انرژی، نرخ گذر از صفر، بسامد پایه، پرچم
واک‌داری، \N{13} ضریب \en{MFCC} و \N{13} دلتای متناظر است.

\subsection{پنجره متن و بردار تصمیم}

تصمیم درباره وقفه بر پایه یک فریم واحد قابل اتکا نیست. به همین دلیل
ویژگی‌های فریم‌های یک «پنجره متن» گرد‌آوری می‌شوند. در این پروژه پنجره
متن \N{400} میلی‌ثانیه است، یعنی \N{20} فریم متوالی. برای هر بعد از بردار
ویژگی فریم، سه آماره میانگین، انحراف معیار و بیشینه روی پنجره محاسبه و در
یک بردار واحد به هم پیوند داده می‌شود. این بردار ورودی طبقه‌بند است.

\section{بازشناسی گفتار}
\label{sec:asr}

\مهم{بازشناسی گفتار} نگاشت یک سیگنال صوتی به رشته‌ای از واژه‌ها است.
معماری‌های امروزی معمولاً یک کدگذار عصبی را با یک هدف آموزشی هم‌ترازی‌آزاد
ترکیب می‌کنند. دو هدف پرکاربرد عبارت‌اند از
«طبقه‌بندی زمانی پیوندگرا»\پاورقی{Connectionist Temporal
Classification (CTC)}~\cite{graves2006ctc} که با فرض استقلال شرطی، احتمال همه
هم‌ترازی‌های ممکن را حاشیه‌ای می‌کند، و
«مبدل عصبی بازگشتی»\پاورقی{RNN Transducer (RNNT)}~\cite{graves2012rnnt}
که یک شبکه پیش‌بین زبانی را نیز در تصمیم دخیل می‌سازد و برای بازشناسی
جویباری مناسب‌تر است.

معماری کدگذار مورد استفاده در سرویس بازشناسی این پروژه از خانواده
\en{Conformer}~\cite{gulati2020conformer} است که لایه‌های
خودتوجهی\پاورقی{self-attention}~\cite{vaswani2017attention} را با لایه‌های
پیچشی ترکیب می‌کند: خودتوجهی وابستگی‌های بلندبرد و پیچش الگوهای محلی طیفی را
مدل می‌کند. مدل مورد استفاده یک مدل ترکیبی \en{RNNT-CTC} با واحدهای زیرواژه
است که در بستر \en{NeMo}~\cite{kuchaiev2019nemo} پیاده‌سازی و به‌صورت یک
سرویس محلی در دسترس قرار گرفته است.

رویکرد جایگزین، آموزش نظارت‌ضعیف در مقیاس بزرگ است؛ نمونه شاخص آن
\en{Whisper}~\cite{radford2023whisper} است که با صدها هزار ساعت داده
چند‌زبانه آموزش دیده و در این پروژه تنها در آزمون‌های مقایسه مؤلفه‌ای به کار
رفته است.

\section{تبدیل متن به گفتار}
\label{sec:tts}

سامانه‌های \مهم{تبدیل متن به گفتار} امروزی معمولاً سرتاسری‌اند. معماری
\en{VITS}~\cite{kim2021vits} که مبنای موتور \en{Piper}~\cite{piper2023} است، یک
خودکدگذار وردشی شرطی را با آموزش هم‌ستیزانه و یک پیش‌بین طول جریان‌محور
ترکیب می‌کند و بدون نیاز به یک نوسان‌ساز جداگانه، مستقیماً موج صوتی تولید
می‌کند.

برای فارسی، مرحله پیش‌پردازش متنی اهمیت خاصی دارد. نوشتار فارسی
واکه‌های کوتاه را ثبت نمی‌کند، فاصله مجازی یا
«نیم‌فاصله»\پاورقی{zero-width non-joiner (ZWNJ)} مرز اجزای واژه‌های
مرکب را تعیین می‌کند، و شکل نوشتاری واژه‌ها ممکن است میان منابع مختلف
متفاوت باشد. به همین دلیل در این پروژه یک «هنجارساز» قطعی متن فارسی
پیش از هر دو مرحله ارزیابی و تولید گفتار اعمال می‌شود تا شکل نویسه‌ها،
نیم‌فاصله‌ها و ارقام یکسان شوند.

استفاده از یک موتور گفتاری «قطعی» با یک مدل صوتی ثابت، یک تصمیم روشمند
است: چون خروجی صوتی برای یک متن ورودی مشخص همیشه یکسان است، اندازه‌گیری‌های
مبتنی بر آن قابل بازتولید می‌شوند.

\section{مدل‌های زبانی بزرگ و تطبیق کم‌پارامتر}
\label{sec:llm}

\subsection{مدل زبانی خودرگرسیو}

یک \مهم{مدل زبانی} احتمال یک رشته نشانه $y_{1:T}$ را به‌صورت زنجیره‌ای
تجزیه می‌کند:

\begin{equation}
p_\theta(y_{1:T}) = \prod_{t=1}^{T} p_\theta(y_t \mid y_{<t}),
\label{eq:lm}
\end{equation}

و با کمینه‌کردن زیان آنتروپی متقاطع روی داده آموزشی برازش می‌شود. معماری
غالب، ترانسفورمر تنها‌کدگشا~\cite{vaswani2017attention} است. مدل‌های
«دستور‌پیرو»\پاورقی{instruction-following} با یک مرحله تنظیم دقیق
نظارت‌شده روی جفت‌های دستور و پاسخ ساخته می‌شوند~\cite{ouyang2022instructgpt}؛
پاسخ‌گوی متنی این پروژه از همین دسته است.

\subsection{تحمیل معلم و اریبی مواجهه}
\label{sec:exposure-bias}

نکته‌ای که برای تفسیر نتایج فصل~\ref{ch:results} تعیین‌کننده است، تفاوت میان
شرایط آموزش و شرایط استنتاج است. در آموزش، مدل در هر گام نشانه‌های
«درست» پیشین را به‌عنوان شرط دریافت می‌کند؛ این شیوه
\مهم{تحمیل معلم} نامیده می‌شود. زیانی که در آموزش و اعتبارسنجی گزارش
می‌شود، تحت همین شرط محاسبه می‌گردد.

اما در استنتاج، مدل باید بر نشانه‌های «تولیدی خودش» شرط بگذارد. اگر
توزیع نشانه‌های تولیدی از توزیع نشانه‌های مرجع فاصله بگیرد، مدل در ناحیه‌ای
از فضای حالت قرار می‌گیرد که در آموزش ندیده است و خطا می‌تواند به‌صورت
تجمعی رشد کند. این ناسازگاری با نام
\مهم{اریبی مواجهه}\پاورقی{exposure bias} شناخته
می‌شود~\cite{ranzato2016scheduledsampling} و راهکار متعارف آن،
«نمونه‌گیری زمان‌بندی‌شده» است که در طول آموزش به‌تدریج نشانه‌های مرجع را
با نشانه‌های تولیدی جانشین می‌کند~\cite{bengio2015scheduledsampling}.

پیامد عملی این مفهوم صریح است: «کاهش زیان تحت تحمیل معلم، به‌خودی‌خود
تضمین نمی‌کند که مدل در تولید خودرگرسیو خروجی قابل استفاده بدهد». آزمایش
مسیر مستقیم این پروژه نمونه‌ای روشن از همین شکاف است.

\subsection{تطبیق کم‌پارامتر و \en{LoRA}}

تنظیم دقیق کامل یک مدل چند‌میلیارد‌پارامتری بر سخت‌افزار محدود عملی نیست.
روش‌های \مهم{تطبیق کم‌پارامتر} تنها شمار کوچکی از پارامترها را آموزش
می‌دهند~\cite{houlsby2019adapters,li2021prefixtuning}. روش مورد استفاده در
این پروژه \en{LoRA}~\cite{hu2022lora} است. ایده آن این است که به‌روزرسانی
مورد نیاز برای یک ماتریس وزن، «رتبه پایین» دارد. برای وزن منجمد
$W_0 \in \mathbb{R}^{d \times k}$، به‌جای یادگیری $\Delta W$ کامل، دو ماتریس
کوچک $A \in \mathbb{R}^{r \times k}$ و $B \in \mathbb{R}^{d \times r}$ با
$r \ll \min(d,k)$ آموزش داده می‌شوند و خروجی لایه چنین می‌شود:

\begin{equation}
h = W_0 x + \frac{\alpha}{r} B A x .
\label{eq:lora}
\end{equation}

در آغاز، $B$ صفر است تا مدل تطبیق‌یافته با مدل پایه یکسان باشد. پارامتر
$r$ رتبه و $\alpha$ ضریب مقیاس است. در این پروژه رتبه‌های \N{64} و \N{128}
آزمایش شده‌اند. مزیت کلیدی \en{LoRA} برای این کار آن است که وزن‌های پایه
منجمد می‌مانند و بنابراین یک آداپتر شکست‌خورده هرگز مدل پایه را خراب
نمی‌کند و می‌توان آن را به‌سادگی برداشت.

\section{رمزینه‌های صوتی عصبی و مدل‌های گفتار‌به‌گفتار}
\label{sec:codec}

\subsection{کوانتیزه‌سازی برداری باقی‌مانده}

برای آن‌که یک مدل زبانی بتواند گفتار تولید کند، صوت باید به رشته‌ای از
نشانه‌های گسسته تبدیل شود. \مهم{رمزینه صوتی عصبی} این کار را با یک
خودکدگذار انجام می‌دهد که در گلوگاه آن یک
«کوانتیزه‌ساز برداری باقی‌مانده»\پاورقی{Residual Vector Quantizer (RVQ)}
قرار دارد. نمونه‌های شاخص \en{SoundStream}~\cite{zeghidour2021soundstream} و
\en{EnCodec}~\cite{defossez2023encodec} هستند.

در کوانتیزه‌سازی باقی‌مانده، بردار نهفته $z$ به‌صورت پی‌درپی تقریب زده
می‌شود. اگر $Q_i$ کتاب رمز $i$‌ام باشد، با $r_0 = z$ داریم:

\begin{equation}
q_i = Q_i(r_{i-1}), \qquad r_i = r_{i-1} - q_i, \qquad
\hat{z} = \sum_{i=1}^{N_q} q_i .
\label{eq:rvq}
\end{equation}

کتاب رمز نخست، اطلاعات درشت و بیشتر «معنایی» را حمل می‌کند و
کتاب‌های رمز بعدی جزئیات «آوایی» را می‌افزایند. این سلسله‌مراتب دلیل
آن است که در آموزش مدل‌های گفتاری، برای زیان کتاب رمز نخست ضریب بزرگ‌تری
در نظر گرفته می‌شود؛ همین ضریب در یکی از آزمایش‌های این پروژه عامل تغییر
بوده است.

\subsection{مدل‌سازی زبانی روی نشانه‌های صوتی}

با گسسته‌شدن صوت، می‌توان همان ماشین مدل‌سازی زبانی را روی نشانه‌های صوتی
اعمال کرد؛ این ایده در \en{AudioLM}~\cite{borsos2023audiolm} و
\en{VALL-E}~\cite{wang2023valle} به کار رفته است. چالش اصلی آن است که هر گام
زمانی نه یک نشانه، بلکه $N_q$ نشانه (یکی برای هر کتاب رمز) دارد. راه‌حل
«ترانسفورمر کوانتیزه‌سازی باقی‌مانده»\پاورقی{RQ-Transformer} است: یک
ترانسفورمر بزرگ «زمانی» روی محور زمان کار می‌کند و یک ترانسفورمر کوچک
«عمقی» در هر گام زمانی، $N_q$ نشانه را به‌صورت خودرگرسیو
تولید می‌کند~\cite{lee2022rqvae}.

\subsection{معماری \en{Moshi} و «تک‌گویی درونی»}
\label{sec:moshi-arch}

مدل \en{Moshi}~\cite{defossez2024moshi} گفت‌وگو را به‌صورت
گفتار‌به‌گفتار مدل می‌کند و سه ویژگی آن برای این پروژه اهمیت دارد.

نخست، \مهم{رمزینه \en{Mimi}}: یک رمزینه صوتی جویباری که بازنمایی معنایی را
با بازنمایی آوایی در کتاب‌های رمز باقی‌مانده تلفیق می‌کند و با گام زمانی
\N{80} میلی‌ثانیه کار می‌کند.

دوم، \مهم{جریان‌های موازی}: مدل به‌طور هم‌زمان دو جریان صوتی را مدل
می‌کند، یکی گفتار خودش و دیگری گفتار کاربر. به همین دلیل مفهوم «نوبت» صریح
حذف می‌شود و هم‌پوشانی، وقفه و بازخورد کلامی به‌طور طبیعی در مدل جای
می‌گیرند. این ویژگی، همان چیزی است که \en{Moshi} را به یک نامزد طبیعی برای
نیازمندی تمام‌دوطرفه بودن تبدیل می‌کند.

سوم، \مهم{تک‌گویی درونی}\پاورقی{Inner Monologue}: پیش از تولید نشانه‌های
صوتی هر گام، مدل نشانه‌های «متنی» هم‌تراز‌شده با زمان را پیش‌بینی
می‌کند. این متن میانی کیفیت زبانی گفتار تولیدی را به‌طور محسوس بهبود
می‌دهد. پیامد آن برای عیب‌یابی این است که می‌توان زیان متنی و زیان صوتی را
جداگانه پایش کرد؛ در فصل~\ref{ch:results} از همین تفکیک برای نشان‌دادن
محل شکست مسیر مستقیم استفاده شده است.

نسخه‌ای از این مدل که در این پروژه به کار رفته \en{Moshika} است؛ یک مدل
هفت‌میلیارد‌پارامتری با صدای زنانه، \N{32} لایه ترانسفورمر، بعد نهفته
\N{4096} و \N{16} کتاب رمز صوتی که هشت مورد از آن‌ها در ترانسفورمر عمقی
پیش‌بینی می‌شوند.

\section{تمام‌دوطرفه بودن و نوبت‌گیری}
\label{sec:duplex}

در گفت‌وگوی انسانی، نوبت‌گیری فرایندی دقیق و سریع است. کار کلاسیک
سازمان نوبت نشان داد که فاصله میان نوبت‌ها کوتاه و منظم است~\cite{sacks1974turntaking}؛
مرورهای مهندسی بعدی، مقدار متعارف این فاصله را در حدود دو‌دهم ثانیه
گزارش کرده‌اند~\cite{skantze2021turntaking}.

از دید مهندسی، سه سطح از قابلیت را باید تفکیک کرد:

\begin{itemize}
\item \مهم{نیم‌دوطرفه}: سامانه در هنگام پخش پاسخ به میکروفون گوش نمی‌دهد.
      کاربر باید منتظر بماند.
\item \مهم{تمام‌دوطرفه با کنترل بیرونی}: میکروفون در هنگام پخش فعال است و یک
      مؤلفه «جداگانه» درباره توقف پخش تصمیم می‌گیرد. این همان
      معماری‌ای است که در این پروژه پیاده و ارزیابی شده است.
\item \مهم{تمام‌دوطرفه ذاتی}: خود مدل مولد، جریان کاربر و جریان خودش را
      هم‌زمان مدل می‌کند و نوبت‌گیری بخشی از توزیع آموخته‌شده است؛ مانند
      \en{Moshi}.
\end{itemize}

در سه رخداد تعاملی که در این نوشتار برچسب‌گذاری می‌شوند تمایز زیر برقرار است:
«وقفه» زمانی است که گوینده دوم آغاز به سخن می‌کند و گفتار او ادامه
می‌یابد و نوبت را می‌گیرد؛ «بازخورد کلامی» یک اظهار کوتاه تأییدی است
که نوبت را نمی‌گیرد؛ و «نوبت پاک» حالتی است که گوینده دوم تنها پس از
پایان گفتار گوینده اول آغاز می‌کند.

\section{طبقه‌بندی با تقویت گرادیانی}
\label{sec:gbdt}

آشکارساز وقفه این پروژه از یک \مهم{مدل تقویت گرادیانی درختی} استفاده
می‌کند~\cite{friedman2001gbm} که در کتابخانه
\en{scikit-learn}~\cite{pedregosa2011sklearn} پیاده‌سازی شده است. این روش یک
مدل جمعی می‌سازد که در آن هر درخت جدید، گرادیان منفی زیان مدل فعلی را
برازش می‌کند:

\begin{equation}
F_m(x) = F_{m-1}(x) + \nu \, h_m(x),
\label{eq:gbm}
\end{equation}

که در آن $h_m$ درخت تازه و $\nu$ نرخ یادگیری است. انتخاب این خانواده مدل
سه دلیل دارد: تعریف پروژه صریحاً یک ماژول مبتنی بر ویژگی‌های دست‌ساز را
خواسته است؛ استنتاج آن روی پردازنده مرکزی در بودجه چند‌میلی‌ثانیه‌ای هر گام
صوتی جای می‌گیرد؛ و با شمار محدود رخدادهای آموزشی، مدل‌های درختی از
شبکه‌های عصبی عمیق پایدارترند.

\section{سنجه‌ها و مفاهیم آماری ارزیابی}
\label{sec:metrics-prelim}

\subsection{نرخ خطای واژه و نویسه}

فاصله \مهم{لِوِنشتاین}~\cite{levenshtein1966} کمینه شمار عملیات درج، حذف و
جانشینی لازم برای تبدیل یک رشته به رشته دیگر است. بر پایه آن،
«نرخ خطای واژه» و «نرخ خطای نویسه» تعریف می‌شوند. در این نوشتار
دو شکل تجمیع گزارش می‌شود. شکل «خرد» که مجموع ویرایش‌ها را بر مجموع
طول مرجع‌ها تقسیم می‌کند:

\begin{equation}
\mathrm{WER}_{\text{micro}} =
\frac{\sum_{i} d\bigl(\hat{w}^{(i)}, w^{(i)}\bigr)}{\sum_{i} \bigl|w^{(i)}\bigr|},
\label{eq:wer}
\end{equation}

و شکل «کلان» که میانگین ساده نرخ خطای هر نمونه است. برای فارسی، تعریف
مرز واژه اهمیت دارد: در این پروژه نشانه‌های واژه با فاصله سفید جدا می‌شوند
و نیم‌فاصله نیز مرز به شمار می‌آید. نشانه‌های نویسه، نقاط کد یونی‌کد بدون
احتساب فاصله سفید و نیم‌فاصله‌اند. همین تعریف باعث می‌شود نرخ خطای واژه در
فارسی بسیار سخت‌گیرانه‌تر از نرخ خطای نویسه باشد؛ نکته‌ای که در
فصل~\ref{ch:discussion} تفسیر می‌شود.

\subsection{سنجه‌های آشکارساز}

برای رخداد دو‌کلاسه «وقفه در برابر غیر‌وقفه»، ماتریس درهم‌ریختگی چهار مقدار
مثبت درست، مثبت نادرست، منفی درست و منفی نادرست را می‌دهد. از آن‌ها
«دقت»، «امتیاز \en{F1}» کلاس وقفه،
«نرخ پذیرش نادرست»\پاورقی{False Accept Rate (FAR)} و
«نرخ رد نادرست»\پاورقی{False Reject Rate (FRR)} محاسبه می‌شوند. در
سامانه‌های تمام‌دوطرفه این دو نرخ خطا هم‌ارز نیستند: پذیرش نادرست باعث
می‌شود پاسخ سامانه بی‌دلیل قطع شود، در حالی که رد نادرست باعث می‌شود
سامانه به وقفه کاربر بی‌اعتنا بماند.

\subsection{فاصله اطمینان و همبستگی گروهی}

برای نسبت‌ها، فاصله اطمینان \مهم{ویلسون}~\cite{wilson1927interval} گزارش
می‌شود که برخلاف تقریب نرمال ساده، برای نمونه‌های کوچک و نسبت‌های نزدیک به
مرز نیز معتبر می‌ماند.

نکته مهم‌تر آن است که رخدادهای برگرفته از یک نشست گفتاری «مستقل»
نیستند؛ گوینده، شرایط آکوستیکی و سبک گفت‌وگو میان آن‌ها مشترک است. بنابراین
یک فاصله اطمینان در سطح رخداد، عدم‌قطعیت واقعی را کم‌برآورد می‌کند. راه‌حل
به‌کاررفته در این پروژه «خودراه‌انداز بلوکی در سطح نشست» است: در هر
تکرار بازنمونه‌گیری، «نشست»ها با جای‌گذاری انتخاب می‌شوند و نه
رخدادهای تک‌تک~\cite{efron1994bootstrap}. در فصل~\ref{ch:results} هر دو
فاصله گزارش شده‌اند تا اثر این تفاوت دیده شود.

\subsection{تفکیک گروهی داده}
\label{sec:group-split}

اگر دو بخش از یک نشست گفتاری واحد در مجموعه آموزش و مجموعه آزمون قرار
بگیرند، مدل می‌تواند گوینده یا موضوع را به‌خاطر بسپارد و نتیجه آزمون
خوش‌بینانه شود. برای پرهیز از این «نشت گروهی»، در این پروژه واحد تفکیک
«شناسه نشست منبع» است: همه جفت‌های برآمده از یک نشست، به‌طور کامل در
یک و تنها یک بخش قرار می‌گیرند. تخصیص با درهم‌سازی قطعی شناسه نشست انجام
می‌شود تا تفکیک بازتولیدپذیر باشد.

\subsection{داوری خودکار با مدل زبانی}
\label{sec:judge}

سنجش کیفیت معنایی پاسخ گفت‌وگو با سنجه‌های سطح رشته ممکن نیست، زیرا برای
یک پرسش، پاسخ‌های درست متعددی وجود دارد. یک جانشین متعارف، استفاده از یک
مدل زبانی بزرگ به‌عنوان «داور» است~\cite{zheng2023mtbench}. در این
روش، پاسخ نامزد و پاسخ‌های مقایسه‌ای به‌صورت «دوسوکور» و با ترتیب
تصادفی‌شده به داور داده می‌شوند و داور بر پایه یک سنجه‌نامه از پیش تعیین‌شده
امتیاز می‌دهد.

محدودیت‌های شناخته‌شده این روش عبارت‌اند از اریبی نسبت به ترتیب ارائه، اریبی
نسبت به طول پاسخ، و اریبی نسبت به پاسخ‌های تولیدشده توسط مدل‌های هم‌خانواده
با داور. در این پروژه سه اقدام برای کاهش این اثرها انجام شده است:
دوسوکورسازی ترتیب، کدگشایی قطعی داور با دمای صفر، و تکرار هر فراخوانی و
گزارش نرخ توافق. با این حال، همان‌طور که در فصل~\ref{ch:discussion} تأکید
شده است، این سنجه یک «تقریب خودکار» است و جانشین ارزیابی انسانی
نمی‌شود؛ به‌ویژه آن‌که در این پروژه داور و نامزد از یک خانواده مدل‌اند.

\subsection{سنجه‌های تأخیر}
\label{sec:latency-prelim}

دو سنجه زمانی در این نوشتار تعریف می‌شوند:

\begin{itemize}
\item $T_{\text{first\_audio}}$: زمان از پایان گفتار کاربر تا نخستین نمونه
      صوتی «شنیدنی» پاسخ. مقدار رسمی این سنجه باید از سمت
      کارخواه گزارش شود، زیرا تنها مرورگر می‌داند که منبع صوت واقعاً
      زمان‌بندی شده است.
\item $T_{\text{barge\_in}}$: زمان از آغاز آوایی گفتار وقفه‌ای کاربر تا
      لحظه‌ای که پخش پاسخ «واقعاً» متوقف می‌شود؛ این مقدار نیز به تأیید
      کارخواه نیاز دارد.
\end{itemize}

نکته روشمند مهم آن است که «زمان تولید کامل پاسخ در سمت سرور»، سنجه
$T_{\text{first\_audio}}$ نیست. زمان تولید سرور یک اندازه‌گیری سودمند برای
عیب‌یابی است، اما نه شرط پذیرش تأخیر. این تفکیک در فصل~\ref{ch:results} با
دقت رعایت شده است.
